\section{Latent Diffusion Models：潜在扩散模型}

\subsection{像素空间扩散的计算瓶颈}

在之前讨论的所有扩散模型中，扩散过程直接在\textbf{像素空间}进行：所有中间变量 $x_t$ 与原始图像具有相同的维度。对于高分辨率图像（如 $512\times512\times3$ 或更大），这带来了巨大的计算和内存开销。

\begin{warningbox}{直接在像素空间训练的代价}
以 Stable Diffusion 为例，其训练数据集 LAION-5B 包含约 58 亿张图像。如果直接在像素空间训练高分辨率扩散模型：
\begin{itemize}
    \item 使用\textbf{单张 GPU}训练需要约 \textbf{2--3 年}
    \item 生成 10,000 张图像（单张 GPU）需要约\textbf{一天}
\end{itemize}
所有中间变量与输出必须保持相同维度，这是扩散模型相比 GAN 的一个结构性劣势。
\end{warningbox}

\subsection{两阶段架构：压缩 + 扩散}

Latent Diffusion Models (LDM) 的核心思想是将扩散过程从像素空间移到\textbf{低维潜在空间}：

\textbf{第一阶段：训练图像自编码器}
\begin{itemize}
    \item 编码器 $\mathcal{E}$：将高分辨率图像 $x \in \mathbb{R}^{H\times W\times 3}$ 压缩为低维潜在表示 $z = \mathcal{E}(x) \in \mathbb{R}^{h\times w\times c}$
    \item 解码器 $\mathcal{D}$：从潜在表示重建图像 $\hat{x} = \mathcal{D}(z) \approx x$
    \item 压缩率：通常 $h = H/8$, $w = W/8$，空间维度缩小 64 倍
\end{itemize}

\textbf{第二阶段：在潜在空间训练扩散模型}
\begin{itemize}
    \item 所有扩散操作（加噪、去噪、score 计算）都在 $z$ 空间进行
    \item 噪声预测网络输入 $z_t$（而非 $x_t$），输出与 $z_t$ 同维度的噪声估计
    \item 生成完成后，通过解码器 $\mathcal{D}$ 将结果映射回像素空间
\end{itemize}

\begin{importantbox}{Latent Diffusion 的关键优势}
\begin{enumerate}
    \item \textbf{计算效率}：在 $64\times64$ 的潜在空间运行扩散，比在 $512\times512$ 的像素空间快约 64 倍
    \item \textbf{语义丰富}：潜在空间通过自编码器训练已编码了高层语义信息，扩散模型可以更关注语义结构而非像素细节
    \item \textbf{模块化}：自编码器和扩散模型可以独立训练和更换
\end{enumerate}
\end{importantbox}

\subsection{自编码器的选择}

LDM 中使用的自编码器并非标准 VAE（其 KL 散度正则化权重较高会牺牲重建质量），而是更接近\textbf{普通自编码器}（autoencoder），重点优化重建质量：

\begin{itemize}
    \item 使用较低权重的 KL 散度正则化，保证潜在空间的正则性
    \item 加入感知损失（perceptual loss）和对抗损失（adversarial loss）提升重建质量
    \item 编码器-解码器通常采用卷积架构，支持任意分辨率输入
\end{itemize}

\begin{knowledgebox}{为什么不用标准 VAE？}
标准 VAE 的 KL 正则化鼓励潜在空间接近标准正态分布，这会损害重建质量。在 LDM 中，自编码器的主要目标是\textbf{无损压缩}——在大幅降低维度的同时尽可能保持图像信息。因此更倾向于使用较弱的正则化，让潜在空间保留更多细节。
\end{knowledgebox}

\subsection{条件注入机制}

在 LDM 的架构中，条件信息（文本、图像、语义图等）通过\textbf{交叉注意力}（cross-attention）机制注入 U-Net：

\begin{enumerate}
    \item 条件信息 $y$（如文本）首先通过条件编码器（如 CLIP text encoder）映射为 embedding 向量
    \item 在 U-Net 的每个 attention block 中，$z_t$ 的特征作为 Query，条件 embedding 作为 Key 和 Value
    \item 交叉注意力的输出融合了 $z_t$ 的空间信息和条件的语义信息
\end{enumerate}

\begin{table}[H]
\centering
\begin{tabular}{lll}
\toprule
\textbf{组件} & \textbf{输入维度} & \textbf{功能} \\
\midrule
图像编码器 $\mathcal{E}$ & $512\times512\times3 \to 64\times64\times4$ & 像素$\to$潜在空间 \\
噪声预测 U-Net & $64\times64\times4 + \text{time} + \text{cond}$ & score 估计 \\
条件编码器（CLIP） & text $\to \mathbb{R}^{77\times768}$ & 文本$\to$embedding \\
图像解码器 $\mathcal{D}$ & $64\times64\times4 \to 512\times512\times3$ & 潜在空间$\to$像素 \\
\bottomrule
\end{tabular}
\caption{Stable Diffusion 核心组件及其维度}
\end{table}

\subsection{Stable Diffusion 架构}

Stable Diffusion 是 LDM 最知名的实现，由 CompVis（LMU Munich）和 Stability AI 合作开发。其完整架构包含：

\begin{enumerate}
    \item \textbf{VAE 编码器/解码器}：8 倍空间下采样，4 通道潜在空间
    \item \textbf{U-Net with cross-attention}：约 860M 参数，包含 ResNet blocks + Transformer blocks
    \item \textbf{CLIP text encoder}：将文本 prompt 编码为 $77\times768$ 的 token embedding 序列
    \item \textbf{CFG 支持}：训练时使用 10\% 的概率将文本条件替换为空字符串
\end{enumerate}

\begin{knowledgebox}{从 U-Net 到 Transformer}
早期 LDM（包括 Stable Diffusion v1/v2）使用 U-Net 作为噪声预测网络。近期趋势是转向纯 Transformer 架构（如 DiT），因其更易扩展、与 NLP 技术栈一致。Stable Diffusion 3 已采用 MM-DiT（多模态 Diffusion Transformer）架构。
\end{knowledgebox}

\subsection{本章小结}

Latent Diffusion Models 通过"先压缩再扩散"的两阶段设计，将高维像素空间的扩散问题转化为低维潜在空间的扩散问题，在大幅降低计算成本的同时保持了生成质量。这一架构已成为现代图像/视频生成模型的标准范式（Stable Diffusion、DALL-E 3、Sora 等），是将扩散模型从学术原型推向工业应用的关键技术创新。

\newpage
